% CS771 progress report -- CVPR 2027 author kit, camera-ready mode, 11pt as required.
% Compile: tectonic progress_report.tex   (or upload this folder to Overleaf)
\documentclass[11pt,twocolumn,letterpaper]{article}
\usepackage[T1]{fontenc}       % makes Times (from cvpr.sty) work under XeTeX/tectonic too
\usepackage{cvpr}              % no [review] option => final / camera-ready version
\usepackage{graphicx}
\usepackage{amsmath}
\usepackage{booktabs}
\usepackage{xcolor}
\definecolor{cvprblue}{rgb}{0.21,0.49,0.74}
\usepackage[breaklinks,colorlinks,allcolors=cvprblue]{hyperref}
\usepackage{enumitem}
\setlist{nosep,leftmargin=*}

\title{Occlusion-Robust ST-GCN for Smart-Environment Action Recognition\\{\large Progress Report}}
\author{Karthik Reddy Jannupalli \and Sanjana Kallingal \and Penchala Siva Prasanth Kasa}

\begin{document}
\maketitle

\begin{table*}[t]
\centering\footnotesize
\setlength{\tabcolsep}{3.5pt}
\begin{tabular}{@{}llcccccccc@{}}
\toprule
 & & \multicolumn{5}{c}{Seen-type occlusion (structured)} & \multicolumn{3}{c}{Held-out occlusion} \\
\cmidrule(lr){3-7}\cmidrule(l){8-10}
Model & Train aug. & Clean & Lower & Furn. & Furn.+noise & Self & Arm$^\dagger$ & Temporal$^\dagger$ & Rand.\,30\%$^\ddagger$ \\
\midrule
Bi-LSTM & -- & 75.9 & 53.9 & 56.3 & 66.9 & 73.0 & 39.8 & 43.6 & 65.1 \\
Bi-LSTM & geometric (ours) & 76.6 & 72.9 & 75.1 & 75.0 & 76.0 & 51.7 & 52.7 & 69.8 \\
1D-CNN & -- & 78.6 & 45.5 & 47.8 & 67.1 & 75.1 & 41.7 & 43.8 & 61.3 \\
1D-CNN & geometric (ours) & 79.1 & 73.7 & 77.4 & 77.1 & 78.4 & 50.9 & \textbf{60.3} & 71.3 \\
\midrule
ST-GCN & -- & \textbf{82.5} & 76.9 & 74.8 & 78.7 & 78.8 & 65.4 & 43.7 & 37.3 \\
ST-GCN & random joint & 82.2 & 76.0 & 74.8 & 79.0 & 80.8 & 65.8 & 36.6 & \textbf{81.4} \\
ST-GCN & lower-body (prop.) & 80.7 & \textbf{80.7} & 75.5 & 78.8 & 77.7 & 63.1 & 43.6 & 42.5 \\
ST-GCN & geometric (ours) & 82.1 & \textbf{80.7} & \textbf{81.2} & \textbf{81.2} & \textbf{81.9} & \textbf{68.8} & 54.8 & 61.0 \\
\midrule
\multicolumn{2}{@{}l}{\emph{Joints hidden (\%)}} & 0 & 32 & 19 & 20 & 3 & 24 & 41 & 30 \\
\bottomrule
\end{tabular}
\caption{Top-1 accuracy (\%) on our 22-class domestic subset of NTU-60 (cross-subject, 6{,}038 test clips, one seed, 30 epochs).
Hidden joints are zeroed except in ``Furniture+noise'', where they hold their last observed position plus jitter.
$^\dagger$Never seen in training by any variant. $^\ddagger$Seen only by the random-joint variant.}
\label{tab:results}
\end{table*}

%--------------------------------------------------------------------
\section{Status and Scope Change}
We have a working end-to-end pipeline and first results (Tab.~\ref{tab:results}).
\textbf{Changes from the proposal} (in response to the feedback, Sec.~\ref{sec:resp}):
(i) the fixed lower-body cluster dropout is replaced by a \emph{geometric} occlusion model covering
furniture and self-occlusion; (ii) the evaluation now uses several occlusion types, including types
never seen in training; (iii) the 2D-keypoint branch (MediaPipe~\cite{lugaresi2019mediapipe}/YOLOv8) and the edge-hardware study are deferred.

%--------------------------------------------------------------------
\section{Work Completed}
\textbf{Data.} 22 domestic single-person classes of NTU RGB+D 60~\cite{shahroudy2016ntu} (\eg drink, read, type, sit down, put on shoe, falling),
cross-subject split (14{,}690 train / 6{,}038 test), Kinect 3D joints in camera coordinates, resampled to 64 frames.
Each sequence is centered on the spine-shoulder joint at its first \emph{visible} frame, so normalization never uses hidden joints.

\textbf{Occlusion model (ours).}
\emph{Furniture:} a virtual box stands on the floor between camera and subject (random height from feet to above the hips, width, offset and distance).
A joint is hidden when its camera ray crosses the box, which yields partial, perspective-correct occlusion that changes per frame as the person moves.
\emph{Self-occlusion:} a joint is hidden when it projects inside the torso (shoulders--hips, widened by body thickness) and lies behind it.
\emph{Fill:} hidden joints are either zeroed or, mimicking sensor behaviour, held at the last observed position plus 5\,cm jitter.
Training applies one of these per sample with probability 0.6, sometimes only on a sub-window of frames.

\textbf{Models.} ST-GCN~\cite{yan2018stgcn} (10 blocks, 3.1M params), Bi-LSTM (2.3M) and 1D-CNN (1.0M) baselines, all trained 30 epochs with AdamW on one NVIDIA L4.
A 69-check test suite verifies the occlusion geometry and the training/eval pipeline.

\textbf{Preliminary findings.}
(1)~Geometric augmentation makes ST-GCN the best or tied-best model on every structured occlusion
(furniture $+6.4$, self-occlusion $+3.1$ over no augmentation) at a 0.4-point clean cost.
It also transfers to held-out types (one arm $+3.4$, temporal block $+11.1$).
(2)~The proposal's lower-body dropout helps \emph{only} on the matching lower-body test ($+3.8$) and barely on realistic furniture occlusion ($+0.7$), which confirms the reviewer's concern.
(3)~Random-joint dropout wins only on random dropout, the condition it trains on, and hurts temporal occlusion.
(4)~Zeroed joints hurt unaugmented models far more than sensor-like noisy fill (1D-CNN 47.8 vs.\ 67.1), so simulation choices strongly change robustness conclusions.
(5)~Whole-frame (temporal) occlusion remains the hardest case for all models ($\leq$60\%).

\textbf{Difficulties.} Official NTU access requires approval, so we used the public preprocessed
release from PYSKL/OpenMMLab~\cite{duan2022pyskl}. The new GCP project had zero GPU quota, which we resolved via a quota request.
All occlusion so far is simulated, and results are from a single seed.

%--------------------------------------------------------------------
\section{Next Steps}
\begin{enumerate}
\item \textbf{Real occlusion.} Evaluate on Toyota Smarthome~\cite{das2019toyota} (real apartments, furniture occlusion; requires an access request) and on NTU clips with untracked Kinect joints (tracking state from the raw data).
\item \textbf{Temporal robustness.} Add frame-block occlusion with temporal interpolation, and feed a per-joint visibility/confidence channel.
\item \textbf{Rigor.} Three seeds, the cross-view split, and the part/frame occlusion protocol of RA-GCN~\cite{song2020ragcn}.
\item \textbf{Deployment.} 2D keypoints from MediaPipe on RGB, and CPU/edge latency. (So far only $\approx$5.4\,ms per clip on a shared L4 GPU.)
\end{enumerate}

%--------------------------------------------------------------------
\section{Response to Proposal Feedback}\label{sec:resp}
\textbf{``Random masking may not simulate realistic self- and object occlusion.''}
We agree. The proposal masked a fixed cluster at random, which does not model how occlusion arises.
We now derive masks from 3D scene geometry (furniture occluders) and from body geometry (self-occlusion), and corrupt hidden joints as a sensor would.
Finding~(2) shows the difference matters: cluster dropout does not transfer to furniture occlusion, but geometric occlusion does.

\textbf{``The evaluation does not consider realistic occlusion.''}
Test-time occlusion now includes furniture, self-occlusion, noisy fill and three held-out types, rather than only the training augmentation.
Real occluded data (Toyota Smarthome, untracked NTU joints) is our first next step.

{\small
\bibliographystyle{ieeenat_fullname}
\bibliography{refs}
}
\end{document}
